% GPT-6 Astra Pro; effort not exposed; /root/replacement; 22 November 2026.
\section[HD-F1-R16: adjacent stable endpoints in virtual rank one]{HD-F1-R16: adjacent stable torsion-free endpoints cannot yield a balanced virtual-rank-one weakly semiregular object}
\label{sec:followup-rm-stable-adjacent}

Use the genuine real-quadratic fourfold $X$, compatible
polarization and whole balanced character of R07--R11. There
is no actual perfect object $E$ with
\[
 \begin{gathered}
 H^0(E)=A,\qquad H^1(E)=T,\qquad H^j(E)=0\ (j\ne0,1),\\
 \operatorname{ch}(E)=N e^{-B}v_q,\qquad
 v_q=(1-qA_1)(1-qA_2),\quad q>0,\quad N\in\{1,-1\},
 \end{gathered}
 \tag{F1R16.1}
\]
whose two positive-rank cohomology sheaves $A,T$ are
torsion free and slope stable of the same slope for one
ample polarization, and which has a weakly semiregular
lift for a coherent faithful finite point action. This
exclusion permits arbitrary singularities of the
torsion-free sheaves and every actual connecting class
in $\operatorname{Ext}^2_X(T,A)$. It requires neither
semihomogeneous endpoints nor simplicity of $E$.

R17 below strengthens this exclusion to every nonzero
virtual rank by adding the independent characteristic
point-action tangents. The present proof is retained
as the intermediate extension-tangent calculation.

We prove the stronger precise exclusion in which stability
and equal slope are replaced by the actual conditions
\[
 \operatorname{End}_X(A)=\operatorname{End}_X(T)=\mathbf C,
 \qquad
 \operatorname{Hom}_X(A,T)=\operatorname{Hom}_X(T,A)=0.
 \tag{F1R16.2}
\]
Both endpoints remain positive-rank torsion-free sheaves.
No existence assertion for a pair satisfying a separate
character equation is used.

\paragraph{1. The stable endpoints give exactly the required Hom vanishings.}
The ranks $r_A,r_T$ are different because
$r_A-r_T=N=\pm1$. For a nonzero morphism from one of
the stable endpoints to the other, let $I$ be its image.
If the kernel has positive rank, stability of the source
gives $\mu(I)>\mu(\text{source})$. A lower-rank image in
the target has $\mu(I)<\mu(\text{target})$ by stability;
a full-rank image has $\mu(I)\le\mu(\text{target})$
because the codimension-one part of the quotient is
effective. Both contradict equality of the endpoint
slopes. If the kernel has rank zero, it is zero because
the source is torsion free; the unequal ranks then make
$I$ a proper lower-rank subsheaf of the target, again
contradicting equality of slopes. This proves both Hom
vanishings in (F1R16.2).

For completeness a slope-stable torsion-free sheaf has
only scalar endomorphisms. The same image-and-kernel
argument makes every nonzero endomorphism have full
generic rank. Its characteristic-polynomial coefficients
are regular on the locally free locus, extend across the
codimension-at-least-two complement on the smooth $X$,
and are constant because $X$ is projective and connected.
Subtracting an eigenvalue times the identity would give
a nonzero endomorphism of smaller generic rank unless
the original endomorphism were scalar. Thus (F1R16.2)
holds. This argument uses the unequal ranks for the
cross-Hom assertion; it makes no assertion that arbitrary
equal-rank stable torsion-free sheaves with the same slope
have no maps between them.

\paragraph{2. The canonical triangle forces both whole endpoint characters.}
Write the actual standard truncation triangle as
\[
 A\xrightarrow{i} E\xrightarrow{p} C:=T[-1]
       \xrightarrow{\delta} A[1],\qquad
       \delta\in\operatorname{Ext}^2_X(T,A).
 \tag{F1R16.3}
\]
The point functors are t-exact, so a specified coherent
lift of $E$ induces coherent lifts of $A,T$ and the whole
triangle. We write $\operatorname{Ext}_G$ for the Ext
groups of these actual lifts. The characteristic action
of $HH^*(X)$ is the one transported by the identity-sector
map of R05; all its ordinary evaluation values are
$G$-invariant.

Put $K_E=\ker c^2_{\operatorname{ch}(E)}$. Weak
semiregularity and the actual forgetful isomorphism on
$HH_{-2}$ imply
$\operatorname{ev}_E(K_E)=0$ in ordinary self-Ext, as
in (F1R11.9). Apply the two endpoint Hom sequences to
(F1R16.3). Their only possible kernels are precisely
the two Yoneda images
\[
 \operatorname{Hom}(A,T)\longrightarrow
 \operatorname{Ext}^2(A,A),\quad f\longmapsto\delta f,
 \qquad
 \operatorname{Hom}(A,T)\longrightarrow
 \operatorname{Ext}^2(T,T),\quad f\longmapsto f[2]\delta,
 \tag{F1R16.4}
\]
proved in R10. They vanish by (F1R16.2). Naturality along
$i,p$ therefore gives
$\operatorname{ev}_A(K_E)=\operatorname{ev}_T(K_E)=0$.
Taking the full trace and applying the complete common
annihilator (F1R11.8) gives the exact whole characters
\[
        \operatorname{ch}(A)=r_Ae^{-B}v_q,
        \qquad
        \operatorname{ch}(T)=r_Te^{-B}v_q.
 \tag{F1R16.5}
\]
The proportionality constants are fixed by the actual
ranks, also when $N=-1$. DEG\_L S-22 excludes a rank-one
torsion-free sheaf with this normalized character. Hence
$r_A,r_T\ge2$, and their difference of absolute value
one gives
\[
                            r_A r_T\ge6.
 \tag{F1R16.6}
\]

\paragraph{3. The endomorphism dimension and actual extension-to-tangent map.}
The four negative-Hom groups between $A$ and $C=T[-1]$
all vanish: the only cross case that is not an automatic
negative sheaf Ext is
$\operatorname{Hom}(C,A[-1])=\operatorname{Hom}(T,A)=0$.
Applying the two long exact Hom sequences to the actual
extension (F1R16.3) consequently gives
$\operatorname{Ext}^j_X(E,E)=0$ for $j<0$. In particular
the lift is gluable.

Set
\[
 \begin{gathered}
 a=\operatorname{ext}^1_G(T,A),\qquad
 b=\operatorname{ext}^1_G(A,T),\qquad
 c=\chi_G(T,A),\\
 V=\operatorname{Ext}^2_G(T,A),\qquad
 \varepsilon=\dim\mathbf C\delta\in\{0,1\},\qquad
 e_j=\operatorname{ext}^j_G(E,E),\quad k=\chi_G(E,E).
 \end{gathered}
 \tag{F1R16.7}
\]
The scalar endomorphisms of each endpoint are invariant,
so their equivariant endomorphism algebras are also
$\mathbf C$. All Hom sequences below may be taken in
the equivariant category; equivalently one can take
invariants, which is exact over $\mathbf C$.

The two groups $\operatorname{Hom}(A,C)$ and
$\operatorname{Hom}(A[1],C)$ vanish. Hence
\[
                 \operatorname{Hom}_G(E,C)=\mathbf C p.
\]
Also $\operatorname{Hom}(E,C[-1])=0$ by standard
t-structure orthogonality. Applying $\operatorname{Hom}(-,A)$
and then $\operatorname{Hom}(E,-)$ to (F1R16.3) gives
\[
 \begin{split}
 0&\longrightarrow\operatorname{Ext}^1_G(T,A)
       \longrightarrow\operatorname{Hom}_G(E,A)
       \longrightarrow\mathbf C
       \xrightarrow{\,1\mapsto\delta\,} V,\\
 0&\longrightarrow\operatorname{Hom}_G(E,A)
       \xrightarrow{\,i\circ(-)\,}\operatorname{End}_G(E)
       \xrightarrow{\,p\circ(-)\,}\mathbf C p
       \longrightarrow0.
 \end{split}
 \tag{F1R16.8}
\]
The last map is onto because it sends $1_E$ to $p$.
Consequently
\[
                             e_0=a+2-\varepsilon.
 \tag{F1R16.9}
\]

The same triangle supplies the actual linear map
\[
 J:V\longrightarrow\operatorname{Ext}^1_G(E,E),
                       \qquad J(v)=i[1]vp.
 \tag{F1R16.10}
\]
Its kernel is exactly $\mathbf C\delta$. Indeed,
$\operatorname{Hom}_G(E,C)=\mathbf C p$ and
$\delta p=0$ make the second arrow in
\[
 \operatorname{Hom}_G(E,C)\xrightarrow{\delta\circ(-)}
 \operatorname{Hom}_G(E,A[1])\xrightarrow{i[1]\circ(-)}
 \operatorname{Ext}^1_G(E,E)
\]
injective. Applying $\operatorname{Hom}(-,A[1])$ to the
triangle shows that the kernel of
$V\xrightarrow{(-)p}\operatorname{Hom}_G(E,A[1])$
is the image of
$\operatorname{End}_G(A[1])\xrightarrow{(-)\delta}V$,
which is $\mathbf C\delta$. Thus
\[
                   V/\mathbf C\delta\lhook\joinrel\longrightarrow
                   \operatorname{Ext}^1_G(E,E),\qquad
                   \dim V-\varepsilon\le e_1.
 \tag{F1R16.11}
\]
This proof includes $\delta=0$; it does not discard the
split triangle or require a nonzero connecting class.

\paragraph{4. Exact Euler scaling gives the contradiction.}
The invariant CY4 pairing of S-21 gives
$\operatorname{Ext}^3_G(T,A)=\operatorname{Ext}^1_G(A,T)^*$.
The cross degree-zero and degree-four groups vanish by
(F1R16.2) and the same pairing. Therefore
\[
                             \dim V=c+a+b.
 \tag{F1R16.12}
\]
Weak semiregularity gives $e_2=18$ by the characteristic
rank and the exact Serre transpose in R07 and R11.
Gluability and invariant CY4 duality then give
\[
 k=2e_0-2e_1+18,
 \qquad
 e_1=a+11-\varepsilon-\frac{k}{2}.
 \tag{F1R16.13}
\]
Combining (F1R16.11)--(F1R16.13) yields
\[
                              c+b+\frac{k}{2}\le11.
 \tag{F1R16.14}
\]

Here the exact cross Euler scaling follows from the
already proved map in S-21, not from a dimension estimate
on invariants. Apply (S21.7) to the actual equivariant
objects $A,T,A\oplus T$. Its proof uses only the
forgetful/induction adjunction, the identity-sector
isomorphism and the categorical Riemann--Roch pairing;
it imposes no condition on their characters or negative
Ext groups. Subtract the first two identities from the
third. Ordinary and invariant CY4 duality identify the
two cross Euler terms, giving
\[
          \chi_G(T,A)=\frac{\chi_X(T,A)}{|G|},
          \qquad
          k=\frac{\chi_X(E,E)}{|G|}.
 \tag{F1R16.15}
\]
Equivalently the underlying actual morphisms are
$F_*I_*=|G|\operatorname{id}$ and
$I_*F_*=|G|\operatorname{id}$ on $HH_0$, followed by
the adjunction $F\dashv I$ and the bilinear Hochschild
Riemann--Roch pairing. These are precisely the maps
proved in S-21 before its specialization to a secant
character.

Since $\operatorname{td}(X)=1$, (F1R16.5) and S-22 give
\[
 \chi_X(T,A)=4Dq^2r_A r_T,\qquad
 \chi_X(E,E)=4Dq^2N^2,
 \qquad c=\frac{r_A r_T}{N^2}\,k,
 \quad D=\int_X\frac{\theta^4}{24}>0.
 \tag{F1R16.16}
\]
Thus $k$ is positive, and (F1R16.13) makes it an even
integer; in particular $k\ge2$. As $N^2=1$, (F1R16.6)
now gives
\[
             c+\frac{k}{2}
             =\left(r_A r_T+\frac12\right)k
             \ge\frac{13}{2}\cdot2=13.
 \tag{F1R16.17}
\]
This contradicts (F1R16.14), since $b\ge0$, and proves
the claimed unconditional exclusion.

\paragraph{Framework, source and exact scope.}
The character identity is the image of the actual
Grothendieck-group relation $[E]=[A]-[T]$ of
(F1R16.3). The endpoint maps are the canonical truncation
arrows and their precise Yoneda kernels, not formal
additivity of self-Ext. The tangent injection is
(F1R16.10), the invariant trace comparison is the
identity-sector forgetful map, and the cross Euler
formula is the polarization of its exact self-Euler
identity. The source foundations are the published
Buchweitz--Flenner characteristic/trace construction,
Perry's finite-action Hochschild decomposition and
Shklyarov's Hochschild Riemann--Roch theorem, with the
exact locators already checked in S-21 and R05--R11.
No unrefereed theorem supplies a premise.

The proof treats actual adjacent full-rank torsion-free
cohomology and all of their couplings in the stated
virtual ranks. For other nonzero virtual ranks the same
calculation yields the necessary inequality
$k(r_A r_T/N^2+1/2)+b\le11$; no general exclusion in
those ranks is inferred here. It also makes no claim
that arbitrary stable endpoints of different slopes
satisfy (F1R16.2), or that arbitrary cohomology endpoints
are stable.
